\begin{textsample}{POS Dim 3 – rr – Score 12.00 – rr/rr\_inf\_6.txt}  \label{ex:f3_pos_156}
3.3 OBJETIVOS DE APRENDIZAGEM E DESENVOLVIMENTO Em razão das especificidades e diferenças desses sujeitos, a aprendizagem e desenvolvimento para a Educação \textbf{Infantil} no currículo do Estado de Roraima conforme os preceitos da BNCC consideram singularidades que pautam a estrutura dos objetivos propostos, no quais, primeira relaciona -se ao fato das \textbf{crianças}, ao longo da Educação \textbf{Infantil}, serem muito dinâmicas, obedecendo a ritmos muito diversos, sendo impossível prever que um determinado objetivo seja alcançado pela maioria das \textbf{crianças} em um mesmo momento. ( BRASIL, 2016 ) Por esse motivo, na BNCC organizam -se os objetivos de aprendizagem, considerando -se três subgrupos etários : - \textbf{Bebês} ( zero a 1 ano e 6 \textbf{meses} ) Os \textbf{bebês} usam a linguagem da não palavra, mas comunicam muitos pensamentos, sensações, expressões, relações, \textbf{desejos} e emoções, dando sinais de extraordinária versatilidade e expressividade aos seus modos de dizer.
Por conta desse enfoque, entendemos que as propostas de \textbf{cuidado} e educação em espaços coletivos devem priorizar a potência de ação dos \textbf{bebês} e das \textbf{crianças} bem \textbf{pequenas}. - \textbf{Crianças} bem \textbf{pequenas} ( 1 ano e 7 \textbf{meses} a 3 anos e 11 \textbf{meses} ) \textbf{Bebês} e \textbf{crianças} bem \textbf{pequenas} utilizam diferentes formas de comunicação e expressão, que \textbf{adultos} precisam observar e escutar com sensibilidade e inteligibilidade : são \textbf{gestos}, expressões faciais, lágrimas, risos, gritos, silêncios, movimentos, \textbf{balbucios}, entre outros modos e formas de estabelecer relações e conexões com o mundo. - \textbf{Crianças} Pequenas ( 4 anos a 5 anos e 11 \textbf{meses} ) A \textbf{criança} é capaz de pronunciar, de forma adequada, todos os fonemas.
Os diálogos e as histórias são \textbf{contados} com detalhes, com noção de tempo e espaço.
Assim, a \textbf{criança} é considerada em sua totalidade garantindo -lhe o respeito as suas necessidades, ao dar -lhe oportunidade de aprendizagem em diferentes situações no seu cotidiano escolar.
A segunda diz respeito a uma perspectiva com relação aos objetivos de aprendizagem e desenvolvimento que não os circunscreve a um único campo disciplinar, mas inclui conhecimentos de naturezas distintas, relevantes para os \textbf{bebês} e as \textbf{crianças} \textbf{pequenas}, relativos às práticas sociais e às linguagens.
Os objetivos propostos procuram fortalecer o compromisso da Educação \textbf{Infantil}, tanto com os direitos das \textbf{crianças} às aprendizagens, quanto com a vivência da infância pela \textbf{criança} nos distintos Campos de Experiências. ( BRASIL, 2016 )

% matched lemmas: adulto, balbucio, bebê, contar, criança, cuidado, desejo, gesto, infantil, mês, pequeno
\end{textsample}
